\documentclass[11pt,a4paper]{article}

\usepackage[margin=1in]{geometry}
\usepackage{amsmath,amssymb}
\usepackage{booktabs}
\usepackage{tabularx}
\usepackage{hyperref}
\usepackage{xcolor}
\usepackage{listings}
\usepackage{graphicx}
\usepackage{caption}
\usepackage{fancyhdr}
\usepackage{microtype}

\hypersetup{
    colorlinks=true,
    linkcolor=teal!80!black,
    citecolor=teal!80!black,
    urlcolor=blue!70!black
}

\pagestyle{fancy}
\fancyhf{}
\rhead{\small 24AI020 | Chirag Jindani}
\lhead{\small MLOps Lab 4--5: Advanced MLflow \& DVC}
\cfoot{\thepage}

\title{\textbf{\Large Lab Assignment 4--5 Report}\\
\large Advanced MLflow Experiment Tracking, Model Registry \& Data Versioning with DVC}
\author{\textbf{Chirag Jindani} \quad (\textbf{Roll No:} 24AI020)\\
Department of Artificial Intelligence\\
Course: Machine Learning Operations (MLOps)}
\date{October 17, 2026}

\begin{document}

\maketitle
\thispagestyle{empty}

\begin{abstract}
In machine learning projects, versioning code with Git is only half the battle. If you change the dataset, tweak hyperparameters, or retrain a model, code alone cannot tell you how a specific prediction came to be. This lab assignment tackles the question: \textit{Can another engineer identify which model is currently running in production, and easily retrieve the exact code and dataset version used to build it?} Using the California Housing dataset, I set up a complete reproducible pipeline combining \textbf{Git} for code and metadata, \textbf{DVC} for content-addressed dataset versioning and remote storage, and \textbf{MLflow} for experiment tracking, model signatures, and Model Registry lifecycle management.
\end{abstract}

\vspace{0.5em}
\hrule
\vspace{1em}

\section{Introduction \& System Overview}
In standard software engineering, Git tracks everything needed to reproduce a build. In machine learning, a deployed model is the result of non-deterministic training over large external data. 

To achieve full reproducibility without making the repository bloated, I divided responsibilities across three complementary tools:
\begin{enumerate}
    \item \textbf{Git}: Tracks source scripts (`src/`), configuration parameters (`configs/config.yaml`), and lightweight text pointers (`.dvc`, `dvc.yaml`, `dvc.lock`).
    \item \textbf{DVC}: Tracks the actual dataset file (`data/dataset.csv`), moving raw data into a local remote storage cache (`../dvc-storage`) indexed by MD5 checksums.
    \item \textbf{MLflow}: Records runs, logs metrics ($MAE, RMSE, R^2$), visual plots, enforces input/output schema signatures, and manages the model lifecycle in the Model Registry (`@production`, `@candidate`, `@archived`).
\end{enumerate}

\section{Part A: Advanced MLflow (Lab 4)}

\subsection{A1. Key Concepts Explained}
\begin{itemize}
    \item \textbf{Model Signature}: A formal definition of what input columns and types a model expects, and what output type it yields. For California Housing, the signature requires 8 numeric columns (\texttt{MedInc}, \texttt{HouseAge}, etc.) of type \texttt{double} and returns a 1D array of \texttt{float64}. If a client request is missing a column or passes strings, MLflow rejects it immediately with an HTTP 400 error before running prediction code.
    \item \textbf{Artifacts}: Files produced during training that belong to a specific run. In this project, each run logged the serialized model pipeline, an evaluation summary JSON, and three plots (residuals, actual vs. predicted, and feature importances).
    \item \textbf{Autologging}: An MLflow utility (\texttt{mlflow.sklearn.autolog()}) that automatically intercepts training and logs common hyperparameters and metrics. While it easily captured Scikit-learn parameters, it did not capture custom audit tags (DVC hashes, Git SHAs) or isolated test diagnostic plots. I used a hybrid approach: autologging for estimator basics, paired with explicit logging for audit metadata.
    \item \textbf{Model Registry}: A centralized model store where trained models are registered under a friendly name (\texttt{HousePriceModel}) with automatic version numbers (v1, v2, v3), tags, and deployment statuses.
    \item \textbf{Version vs. Stage / Alias}: A \textit{version} is an immutable checkpoint (Version 1, Version 2) that never changes. A \textit{stage} (\texttt{Staging}, \texttt{Production}, \texttt{Archived}) or \textit{alias} (\texttt{@production}, \texttt{@candidate}) is a movable pointer indicating where that version is currently deployed.
\end{itemize}

\subsection{A3. Candidate Model Tracking and Comparison}
I trained three candidate models to evaluate different trade-offs:
\begin{enumerate}
    \item \textbf{Candidate 1 (\texttt{v1\_ridge})}: Regularized linear regression ($\alpha = 1.0$) with feature scaling. Simple, fast, but high bias.
    \item \textbf{Candidate 2 (\texttt{v2\_random\_forest})}: Random Forest Regressor (100 trees, maximum depth 10) capturing non-linear interactions.
    \item \textbf{Candidate 3 (\texttt{v3\_gradient\_boosting})}: Gradient Boosting Regressor (200 trees, learning rate 0.08, depth 5) minimizing residual variance.
\end{enumerate}

\begin{figure}[htbp]
\centering
\fbox{\parbox{0.9\textwidth}{\centering \textbf{[SCREENSHOT 1: MLflow UI Experiment Comparison Table]}\\
\small Open \texttt{http://127.0.0.1:5000}. Capture Experiment \texttt{California\_House\_Price\_Prediction} showing runs \texttt{candidate\_v1\_ridge}, \texttt{candidate\_v2\_random\_forest}, and \texttt{candidate\_v3\_gradient\_boosting} with metrics and Git/DVC tags.}}
\caption{MLflow Tracking UI comparing candidate training runs and metrics.}
\end{figure}

\subsection{A4. Model Signatures and Input Examples}
Each candidate model was logged with an explicit signature inferred via \texttt{infer\_signature(X\_train, y\_pred)}:
\begin{itemize}
    \item \textbf{Inputs}: 8 features of type \texttt{double} (\texttt{MedInc}, \texttt{HouseAge}, \texttt{AveRooms}, \texttt{AveBedrms}, \texttt{Population}, \texttt{AveOccup}, \texttt{Latitude}, \texttt{Longitude}).
    \item \textbf{Outputs}: 1D array of \texttt{float64} representing predicted house values in \$100,000 units.
\end{itemize}

\paragraph{Discussion on Bad Inputs:}
If an API client sends only 3 columns instead of 8, Scikit-learn normally throws a low-level error. Worse, if unlabelled arrays are passed, columns could get silently swapped! With an MLflow signature, incoming data is checked up front. If a column is missing or types mismatch, it returns an explicit HTTP 400 error. 

However, signatures \textit{cannot guarantee} protection against semantic errors: if someone passes a negative house age ($-50$) or coordinates in another state, the signature accepts it because the types are valid floats. Catching out-of-range values requires runtime validation tools.

\begin{figure}[htbp]
\centering
\fbox{\parbox{0.9\textwidth}{\centering \textbf{[SCREENSHOT 2: MLflow Artifacts \& Model Signature View]}\\
\small In MLflow UI, open \texttt{candidate\_v3\_gradient\_boosting}. Select the \texttt{model} directory in Artifacts, displaying the Inferred Input/Output Schema and Input Example.}}
\caption{Inferred model signature and input example recorded in MLflow artifacts.}
\end{figure}

\subsection{A5 \& A6. Model Registry and Lifecycle Decisions}
All three candidates were registered under the unified model name: \texttt{HousePriceModel}. 

An automated acceptance rule was defined in \texttt{configs/config.yaml}:
\[
\text{Acceptance Rule:} \quad \text{Test RMSE} \le 0.65 \quad \text{and} \quad R^2 \ge 0.70
\]

\begin{table}[htbp]
\centering
\small
\caption{Model Registry Versions and Lifecycle Decisions}
\begin{tabularx}{\textwidth}{c c c c c X l}
\toprule
\textbf{Version} & \textbf{Source Run ID} & \textbf{RMSE} & \textbf{MAE} & \textbf{$R^2$} & \textbf{Model Description} & \textbf{Lifecycle Decision} \\
\midrule
\textbf{v1} & \texttt{1d0d1063} & 0.7530 & 0.5315 & 0.5715 & Ridge Baseline. Fails RMSE threshold. & \textbf{Archived / Retired} \\
\textbf{v2} & \texttt{6f1b24b2} & 0.5428 & 0.3639 & 0.7774 & Random Forest (100 trees, depth 10). & \textbf{Staging / Candidate} \\
\textbf{v3} & \texttt{d5f0fe90} & \textbf{0.4783} & \textbf{0.3175} & \textbf{0.8272} & Gradient Boosting (200 trees, lr 0.08). & \textbf{Production} \\
\bottomrule
\end{tabularx}
\end{table}

\paragraph{Critical Thinking: Should the highest-accuracy model always be promoted?}
No. Picking purely based on test RMSE is risky in production:
\begin{enumerate}
    \item \textbf{Inference Latency}: A complex ensemble might take 50ms per prediction, while a simpler model takes 2ms, violating real-time API SLAs.
    \item \textbf{Subgroup Fairness}: Aggregate accuracy can hide severe errors on specific lower-income or rural census blocks.
    \item \textbf{Overfitting}: Flexible models can overfit noise in the test split that won't hold under real-world distribution drift.
    \item \textbf{Asymmetric Error Costs}: In appraisal, underestimating property value carries different financial and customer friction than an overestimation.
\end{enumerate}

\begin{figure}[htbp]
\centering
\fbox{\parbox{0.9\textwidth}{\centering \textbf{[SCREENSHOT 3: MLflow Model Registry Versions Overview]}\\
\small Open the "Models" tab in MLflow UI. Capture the table showing \texttt{HousePriceModel} with Version 1 (Archived), Version 2 (Staging), and Version 3 (Production).}}
\caption{Model Registry versions and assigned operational lifecycle states.}
\end{figure}

\section{Part B: Data Versioning with DVC (Lab 5)}

\subsection{B1. Why DVC Complements Git}
Git is designed for text diffs, not big data files. Committing large CSVs directly into Git creates a new copy in `.git/objects` on every edit, quickly blowing up repo size and hitting GitHub's 100MB limit. DVC solves this by storing the large files in an external cache/remote, while committing a tiny 4-line text pointer file (`.dvc`) containing an MD5 hash into Git.

\subsection{B5 \& B6. Dataset Version Evolution and Recovery Verification}
Two dataset versions were tracked and verified:
\begin{itemize}
    \item \textbf{Dataset v1 (Baseline)}: 20,640 records (1,895,154 bytes, MD5: \texttt{ca5b757f0729c68d74d05bf9d9ca8486}). Committed to Git at \texttt{6767005}.
    \item \textbf{Dataset v2 (Evolution)}: 20,740 records (1,903,729 bytes, MD5: \texttt{5968e75c88d31ad1a1083b00d57fba4e}) containing 100 new survey samples. Committed to Git at \texttt{666c6cd}.
\end{itemize}

\begin{table}[htbp]
\centering
\small
\caption{Dataset Version Comparison and Restoration Verification}
\begin{tabularx}{\textwidth}{c c c c X c}
\toprule
\textbf{Version} & \textbf{Git Commit} & \textbf{Rows} & \textbf{Cols} & \textbf{Change Description} & \textbf{Restoration Verified?} \\
\midrule
\textbf{Dataset v1} & \texttt{6767005} & 20,640 & 9 & Baseline California Housing dataset. & \textbf{Yes} (Verified via \texttt{dvc checkout}) \\
\textbf{Dataset v2} & \texttt{666c6cd} & 20,740 & 9 & Appended 100 new survey records with jitter. & \textbf{Yes} (Verified via \texttt{dvc checkout}) \\
\bottomrule
\end{tabularx}
\end{table}

\paragraph{Verification Method:}
Restoration was verified by checking out the v1 commit on a test branch, running \texttt{dvc checkout}, and running \texttt{wc -l data/dataset.csv}. It restored exactly 20,641 lines (1 header + 20,640 rows) with the matching MD5 hash. A filename alone (\texttt{dataset.csv}) is not enough because file contents can change without the name changing. Only content hashes guarantee that data has not mutated.

\subsection{B7. Reproducible DVC Pipelines}
I created a pipeline stage in \texttt{dvc.yaml} connecting \texttt{data/dataset.csv} to the training script. Running \texttt{dvc repro} executes training and logs hashes to \texttt{dvc.lock}. Running it a second time skips execution because no dependencies changed.

\section{Part C: Integrated MLOps Workflow}

\subsection{C1. Tool Responsibilities Matrix}

\begin{table}[htbp]
\centering
\small
\caption{Comparative Overview: Git, DVC, and MLflow}
\begin{tabularx}{\textwidth}{l X X X}
\toprule
\textbf{Tool} & \textbf{What it Tracks} & \textbf{Storage Location} & \textbf{Why It Matters} \\
\midrule
\textbf{Git} & Code, configs (\texttt{config.yaml}), `.dvc` pointer files, and pipeline manifests (\texttt{dvc.yaml}). & Local repository (\texttt{.git/}) and GitHub. & Tracks exact code history and binds data pointers to pipeline scripts. \\
\textbf{DVC} & Raw datasets, cached blobs, and dependency DAGs. & Local cache (\texttt{.dvc/cache}) and remote (\texttt{../dvc-storage}). & Avoids repo bloat while allowing exact data rollbacks. \\
\textbf{MLflow} & Experiments, hyperparameters, evaluation metrics ($MAE, RMSE, R^2$), visual plots, and Model Registry stages. & Database (\texttt{sqlite:///mlflow.db}) and artifact store (\texttt{mlartifacts/}). & Audits training runs, logs signatures, and manages the production model. \\
\bottomrule
\end{tabularx}
\end{table}

\subsection{C2. Anomaly Investigation Walkthrough}
When an anomalous prediction occurs in production:
\begin{enumerate}
    \item \textbf{Identify Production Model}: Load \texttt{models:/HousePriceModel@production}, revealing Version 3 (Run ID \texttt{d5f0fe90}).
    \item \textbf{Check Audit Tags}: Read \texttt{git\_commit\_sha} (\texttt{666c6cd8}) and \texttt{dvc\_dataset\_hash} (\texttt{5968e75c}).
    \item \textbf{Restore Code}: Run \texttt{git checkout 666c6cd8}.
    \item \textbf{Restore Dataset}: Run \texttt{dvc checkout}, reconstituting the exact 20,740 rows.
    \item \textbf{Debug Locally}: Run \texttt{python src/evaluate.py --model-uri ...} to reproduce the exact environment and weights.
\end{enumerate}

\begin{figure}[htbp]
\centering
\fbox{\parbox{0.9\textwidth}{\centering \textbf{[SCREENSHOT 4: MLflow Run Tags Displaying Git and DVC Audit Hashes]}\\
\small In MLflow UI, expand the Tags section of Run \texttt{candidate\_v3\_gradient\_boosting}. Display \texttt{git\_commit\_sha}, \texttt{dvc\_dataset\_hash}, \texttt{model\_variant}, and \texttt{dataset\_rows}.}}
\caption{Audit tags binding Git commit SHA and DVC data hash to MLflow run.}
\end{figure}

\section{Challenges Encountered and Fixes}
\begin{itemize}
    \item \textbf{Python 3.15 Wheel Issues}: Initial setup with Python 3.15 failed because pre-compiled wheels for \texttt{pygit2} were unavailable on PyPI. I switched the virtual environment to Python 3.14, where wheels installed cleanly.
    \item \textbf{MLflow 3.x Tree Serialization Warning}: Saving tree models threw an untrusted type error with \texttt{skops}. I resolved this by specifying \texttt{serialization\_format="cloudpickle"}.
    \item \textbf{DVC Pipeline Output Conflict}: Declaring \texttt{data/dataset.csv} as both a `.dvc` output and a pipeline output caused an error. I resolved this by setting the dataset as a dependency of the training stage rather than an output.
\end{itemize}

\section{Conclusion}
This lab showed that true reproducibility requires versioning code, data, and model metadata together. By embedding Git commit hashes and DVC content checksums into MLflow Model Registry tags, any production model can be audited, debugged, and restored in minutes.

\end{document}
